\begin{references}
\item Limmathurotsakul D, Turner EL, Wuthiekanun V, et al. Fool’s gold: why imperfect reference tests are undermining the evaluation of novel diagnostics: a reevaluation of 5 diagnostic tests for leptospirosis. Clin Infect Dis. 2012;55(3):322-331.
\item Newman T, Kohn M. Evidence-Based Diagnosis: An Introduction to Clinical Epidemiology. 2nd ed. Cambridge University Press; 2020:89-91.
\item Bland JM, Altman DG. Statistical methods for assessing agreement between two methods of clinical measurement. Lancet. 1986;1(8476):307-310.
\item Newman T, Kohn M. Evidence-Based Diagnosis: An Introduction to Clinical Epidemiology. 2nd ed. Cambridge University Press; 2020:110-137.
\item Tokuda Y, Miyasato H, Stein GH, Kishaba T. The degree of chills for risk of bacteremia in acute febrile illness. Am J Med. 2005;118(12):1417.
\item Sawaya GF, Washington AE. Cervical cancer screening: which techniques should be used and why? Clin Obstet Gynecol. 1999;42(4):922-938.
\item Bjornevik K, Munger KL, Cortese M, et al. Serum neurofilament light chain levels in patients with presymptomatic multiple sclerosis. JAMA Neurol. 2020;77(1):58-64.
\item Smith-Bindman R, Chu P, Miglioretti DL, et al. Physician predictors of mammographic accuracy. J Natl Cancer Inst. 2005;97(5):358-367.
\item Rocker G, Cook D, Sjokvist P, et al. Clinician predictions of intensive care unit mortality. Crit Care Med. 2004;32(5):1149-1154.
\item Newman TB, Puopolo KM, Wi S, Draper D, Escobar GJ. Interpreting complete blood counts soon after birth in newborns at risk for sepsis. Pediatrics. 2010;126(5):903-909.
\sourcepage{446}
\item Newman T, Kohn M. Evidence-Based Diagnosis: An Introduction to Clinical Epidemiology. 2nd ed. Cambridge University Press; 2020:16-22.
\item Newman T, Kohn M. Evidence-Based Diagnosis: An Introduction to Clinical Epidemiology. 2nd ed. Cambridge University Press; 2020:144-167.
\item Moons KG, Altman DG, Reitsma JB, et al. Transparent reporting of a multivariable prediction model for Individual Prognosis or Diagnosis (TRIPOD): explanation and elaboration. Ann Intern Med. 2015;162(1):W1-W73.
\item Moons KGM, Wolff RF, Riley RD, et al. PROBAST: a tool to assess risk of bias and applicability of prediction model studies: explanation and elaboration. Ann Intern Med. 2019;170(1):W1-W33.
\item Food and Drug Administration. Software as a Medical Device (SaMD). \url{https://www.fda.gov/medical-devices/digital-health-center-excellence/software-medical-device-samd}
\item Goff DC Jr, Lloyd-Jones DM, Bennett G, et al. 2013 ACC/AHA guideline on the assessment of cardiovascular risk: a report of the American College of Cardiology/American Heart Association Task Force on Practice Guidelines. Circulation. 2014;129(25 Suppl 2):S49-S73.
\item Grady D, Berkowitz SA. Why is a good clinical prediction rule so hard to find? Arch Intern Med. 2011;171(19):1701-1702.
\item Puopolo KM, Draper D, Wi S, et al. Estimating the probability of neonatal early-onset infection on the basis of maternal risk factors. Pediatrics. 2011;128(5):e1155-e1163.
\item Escobar GJ, Puopolo KM, Wi S, et al. Stratification of risk of early-onset sepsis in newborns \ensuremath{\geq}34 weeks’ gestation. Pediatrics. 2014;133(1):30-36.
\item Kuzniewicz MW, Walsh EM, Li S, Fischer A, Escobar GJ. Development and implementation of an early-onset sepsis calculator to guide antibiotic management in late preterm and term neonates. Jt Comm J Qual Patient Saf. 2016;42(5):232-239.
\item Kuppermann N, Dayan PS, Levine DA, et al. A clinical prediction rule to identify febrile infants 60 days and younger at low risk for serious bacterial infections. JAMA Pediatr. 2019;173(4):342-351.
\item Newman T, Kohn M. Evidence-Based Diagnosis: An Introduction to Clinical Epidemiology. 2nd ed. Cambridge University Press; 2020:175-200.
\item Robinson WR, Renson A, Naimi AI. Teaching yourself about structural racism will improve your machine learning. Biostatistics. 2020;21(2):339-344.
\item Steyerberg EW, Bleeker SE, Moll HA, Grobbee DE, Moons KG. Internal and external validation of predictive models: a simulation study of bias and precision in small samples. J Clin Epidemiol. 2003;56(5):441-447.
\sourcepage{447}
\item Steyerberg EW, Harrell FE Jr. Prediction models need appropriate internal, internal-external, and external validation. J Clin Epidemiol. 2016;69:245-247.
\item Tarnoki DL, Tarnoki AD, Richter A, Karlinger K, Berczi V, Pickuth D. Clinical value of whole-body magnetic resonance imaging in health screening of general adult population. Radiol Oncol. 2015;49(1):10-16.
\item Siegel DL, Edelstein PH, Nachamkin I. Inappropriate testing for diarrheal diseases in the hospital. JAMA. 1990;263(7):979-982.
\item Hitchcock MM, Gomez CA, Banaei N. Low yield of FilmArray GI panel in hospitalized patients with diarrhea: an opportunity for diagnostic stewardship intervention. J Clin Microbiol. 2018;56(3).
\item Lam SHF, Sivitz A, Alade K, et al. Comparison of ultrasound guidance vs. clinical assessment alone for management of pediatric skin and soft tissue infections. J Emerg Med. 2018;55(5):693-701.
\item Viaux-Savelon S, Dommergues M, Rosenblum O, et al. Prenatal ultrasound screening: false positive soft markers may alter maternal representations and mother-infant interaction. PLoS One. 2012;7(1):e30935.
\item Thulin T, Hammar U, Ekbom A, Hultcrantz R, Forsberg AM. Perforations and bleeding in a population-based cohort of all registered colonoscopies in Sweden from 2001 to 2013. United European Gastroenterol J. 2019;7(1):130-137.
\item Zapka J, Taplin SH, Price RA, Cranos C, Yabroff R. Factors in quality care— the case of follow-up to abnormal cancer screening tests—problems in the steps and interfaces of care. J Natl Cancer Inst Monogr. 2010;2010(40):58-71.
\item Welch HG, Schwartz LM, Woloshin S. Overdiagnosed: Making People Sick in Pursuit of Health. Beacon Press; 2011.
\item Esserman LJ, Thompson IM, Reid B, et al. Addressing overdiagnosis and overtreatment in cancer: a prescription for change. Lancet Oncol. 2014;15(6):e234-e242.
\item Detsky AS. A piece of my mind. Underestimating the value of reassurance. JAMA. 2012;307(10):1035-1036.
\item Pletcher MJ, Pignone M. Evaluating the clinical utility of a biomarker: a review of methods for estimating health impact. Circulation. 2011;123(10):1116-1124.
\item Pletcher MJ, Pignone M, Earnshaw S, et al. Using the coronary artery calcium score to guide statin therapy: a cost-effectiveness analysis. Circ Cardiovasc Qual Outcomes. 2014;7(2):276-284.
\item Sheline Y, Kehr C. Cost and utility of routine admission laboratory testing for psychiatric inpatients. Gen Hosp Psychiatry. 1990;12(5):329-334.
\sourcepage{448}
\item Turnbull TL, Vanden Hoek TL, Howes DS, Eisner RF. Utility of laboratory studies in the emergency department patient with a new-onset seizure. Ann Emerg Med. 1990;19(4):373-377.
\item Kohn MA, Carpenter CR, Newman TB. Understanding the direction of bias in studies of diagnostic test accuracy. Acad Emerg Med. 2013;20:1194-1206.
\item Welch HG, Kramer BS, Black WC. Epidemiologic signatures in cancer. N Engl J Med. 2019;381(14):1378-1386.
\item Bell KJ, Del Mar C, Wright G, Dickinson J, Glasziou P. Prevalence of incidental prostate cancer: a systematic review of autopsy studies. Int J Cancer. 2015;137(7):1749-1757.
\end{references}
